% !TEX root =./ECE448_Practice_main.tex
%
% L18 practice set (ADS-B). Every question is a step in the L18 flowgraphs:
% reading PPM chips, splitting a message into its fields, choosing the receiver's
% sample rate and time sink, and explaining what noise does to the table.
% Blank copy: \begin{solution}[h] reserves h of work space.
% SOLUTIONS copy (\iskey=1) prints the red worked solutions.

\fullwidth{\textbf{Learning Objective 6.7: } I can receive and decode ADS-B messages from
aircraft transponders.
\\[4pt]\normalfont\small Show your work. Write a \textbf{Documentation} line at the
end (\textbf{None} if you did not collaborate).}

\begin{questions}

%%%%%%%%%%%%%%%%%%%%%%%%%%%%%%%%%%%%%%%%%
\question \textbf{Reading the chips.} A receiver samples a message at 2 MS/s, so each
0.5~\textmu s half of a bit is one sample, or chip. After the preamble, the first 16
chips it reads, in pairs, are
\begin{center}\texttt{10 01 01 01 10 10 01 10}\end{center}
\begin{parts}

	\begin{minipage}[t]{0.58\linewidth}
		\part Write the 8 bits these chips carry, and the byte in hex.
		\begin{solution}[1.1in]
		In PPM a pulse in the first half (\texttt{10}) is a 1 and a pulse in the second
		half (\texttt{01}) is a 0, so the bits are \texttt{1000 1101}, which is 8D in hex.
		\end{solution}
	\end{minipage}\hfill%
	\begin{minipage}[t]{0.37\linewidth}
		\raggedleft \vspace{12pt}
		bits: \ansbox[1.2in]{1000 1101}

		\vspace{4pt}
		hex: \ansbox[1.0in]{8D}
	\end{minipage}

	\begin{minipage}[t]{0.58\linewidth}
		\part The first 5 bits are the Downlink Format and the next 3 the capability.
		Find both, and say whether this is an ADS-B extended squitter.
		\begin{solution}[0.9in]
		DF $= 10001_2 = 17$ and CA $= 101_2 = 5$. DF 17 is the extended squitter, so yes.
		\end{solution}
	\end{minipage}\hfill%
	\begin{minipage}[t]{0.37\linewidth}
		\raggedleft \vspace{12pt}
		DF: \ansbox[0.8in]{17}

		\vspace{4pt}
		CA: \ansbox[0.8in]{5}
	\end{minipage}

	\part A noise spike makes the sixth pair read \texttt{11}: both chips high. In words,
	explain what the receiver does with that bit, and why the parity check matters here.
		\begin{solution}[1.2in]
		The receiver compares the two chips and calls the bit by whichever is larger, so it
		guesses: with both high it may call the bit wrong. It cannot tell from the chips
		alone. The 24 parity bits catch the wrong bit: the remainder no longer matches, and
		the decoder drops the whole message rather than report a wrong altitude or address.
		\end{solution}

\end{parts}

\newpage
%%%%%%%%%%%%%%%%%%%%%%%%%%%%%%%%%%%%%%%%%
\question \textbf{The message.} The Mode S Decoder prints
\begin{center}\texttt{8DAE07C220483234D71CA031ACD8}\end{center}
The 28 hex digits are the 112 bits: DF and CA (8 bits), ICAO (24), ME (56), and PI (24).
\begin{parts}

	\begin{minipage}[t]{0.58\linewidth}
		\part Split the message into its ICAO address, ME field, and parity, in hex.
		\begin{solution}[1.0in]
		After the first byte, 8D: ICAO is the next 6 hex digits, AE07C2; ME is the next
		14, 20483234D71CA0; PI is the last 6, 31ACD8.
		\end{solution}
	\end{minipage}\hfill%
	\begin{minipage}[t]{0.37\linewidth}
		\raggedleft \vspace{12pt}
		ICAO: \ansbox[1.2in]{AE07C2}

		\vspace{4pt}
		PI: \ansbox[1.2in]{31ACD8}
	\end{minipage}

	\begin{minipage}[t]{0.58\linewidth}
		\part The type code is the first 5 bits of ME. Find it, and say which kind of
		message this is.
		\begin{solution}[0.9in]
		ME begins with 20 hex, \texttt{0010 0000}, so the type code is $00100_2 = 4$:
		an identification message, which carries a callsign.
		\end{solution}
	\end{minipage}\hfill%
	\begin{minipage}[t]{0.37\linewidth}
		\raggedleft \vspace{12pt}
		\ansbox[1.2in]{4, identification}
	\end{minipage}

	\begin{minipage}[t]{0.58\linewidth}
		\part After the type code and 3 category bits, the callsign is eight 6-bit
		characters. In the Mode S character set, 1 to 26 are A to Z, 48 to 57 are the
		digits 0 to 9, and 32 is a space. The eight values are 18, 3, 8, 52, 53, 49, 50, 32.
		Write the callsign.
		\begin{solution}[0.9in]
		18 is R, 3 is C, 8 is H, and 52, 53, 49, 50 are the digits 4, 5, 1, 2; 32 is a
		space. The callsign is RCH4512, a USAF airlift flight.
		\end{solution}
	\end{minipage}\hfill%
	\begin{minipage}[t]{0.37\linewidth}
		\raggedleft \vspace{12pt}
		\ansbox[1.2in]{RCH4512}
	\end{minipage}

\end{parts}

\newpage
%%%%%%%%%%%%%%%%%%%%%%%%%%%%%%%%%%%%%%%%%
\question \textbf{Receiver settings.} The L18 receiver runs at 2 MS/s.
\begin{parts}

	\begin{minipage}[t]{0.58\linewidth}
		\part How many samples does one whole message take, preamble included? Set the
		triggered time sink's number of points to show one whole message with about
		10~\textmu s to spare.
		\begin{solution}[1.0in]
		A message is $8 + 112 = 120$~\textmu s, and at 2 MS/s that is 240 samples. With
		10~\textmu s (20 samples) to spare, about 260 points.
		\end{solution}
	\end{minipage}\hfill%
	\begin{minipage}[t]{0.37\linewidth}
		\raggedleft \vspace{12pt}
		samples: \ansbox[0.9in]{240}

		\vspace{4pt}
		points: \ansbox[0.9in]{about 260}
	\end{minipage}

	\part A classmate sets the RTL-SDR to 1 MS/s to save processing. Explain, in words,
	why the Mode S Decoder can no longer read the bits.
		\begin{solution}[1.1in]
		At 1 MS/s there is one sample per bit, not one per half bit, so the receiver can
		no longer compare the first half of each bit with the second. Each sample may land
		on the pulse or on the quiet half, and nothing tells it which half held the pulse.
		PPM needs at least two samples a bit.
		\end{solution}

	\part On a live run, the 20~ms plot shows a noise floor near 0.05 and messages
	whose pulses reach between 0.15 and 0.6. The trigger level is 0.3. Choose a better
	trigger level and say what you would see with each setting.
		\begin{solution}[1.2in]
		At 0.3 the plot freezes only on the stronger aircraft, whose pulses reach above it;
		the weaker ones never trigger it. A level near 0.1, above the floor but below the
		weakest pulses, triggers on every message. Much lower and the noise itself
		triggers it, and the plot shows noise instead of a preamble.
		\end{solution}

\end{parts}

\newpage
%%%%%%%%%%%%%%%%%%%%%%%%%%%%%%%%%%%%%%%%%
\question \textbf{What noise does.} You run \texttt{L18-adsb-rx.grc} on the simulated
recording and raise the Added noise slider a step at a time.
\begin{parts}

	\part Describe what happens to the 20~ms plot, the Messages label, and the aircraft
	table as the noise rises.
		\begin{solution}[1.2in]
		The noise floor in the plot rises toward the spikes, and the weaker spikes sink
		into it. The count of messages that pass the parity drops, and the table starts
		over and fills with fewer aircraft: the strongest ones stay and the weakest ones
		never return.
		\end{solution}

	\part Explain, in words, why the weakest aircraft drop out first, and what that
	means for a real receiver.
		\begin{solution}[1.2in]
		A weak aircraft's pulses stand only a little above the noise, so a noise peak in
		the quiet half of a bit beats the pulse in some pair of almost every message, the
		parity fails, and the message is dropped. A strong aircraft's pulses are far above
		the noise, so its bits survive. On a real receiver the weak aircraft are the
		distant ones, so the range at which aircraft decode shrinks as the noise rises.
		\end{solution}

\end{parts}

\newpage
%%%%%%%%%%%%%%%%%%%%%%%%%%%%%%%%%%%%%%%%%
\question \textbf{Generating and locating.} Anyone with an RTL-SDR can receive ADS-B, and a flowgraph can generate it.
\begin{parts}

	\part The L18 generator writes a file. Explain why nothing in the course transmits
	its output at 1090 MHz.
		\begin{solution}[1.1in]
		Air traffic control and aircraft collision-avoidance systems use 1090 MHz, and
		nothing in an ADS-B message proves who sent it, so a transmitted message could
		put a false aircraft on a controller's display or another aircraft's. Transmitting
		there is illegal without authorization and unsafe.
		\end{solution}

	\part The Mode S Decoder finds each position from one message by using USAFA as a
	reference position. Explain, in words, why it needs a reference at all, and what goes
	wrong for an aircraft 500 nautical miles away.
		\begin{solution}[1.2in]
		A position message carries the latitude and longitude only within one of many zones
		(Compact Position Reporting), not which zone. The decoder picks the zone nearest its
		reference, which is right for any aircraft within about 180 nautical miles. An
		aircraft 500 nautical miles away lies in a different zone, so the decoder places it
		in the wrong one and reports a position that is badly wrong; pairing an even and an
		odd message avoids that.
		\end{solution}

\end{parts}

\end{questions}
